\subsection{紧自同态的对角化}

定理 1. --- 假设 K = C。设 u 是 E 的紧且正规的自同态。存在 E 的一个标准正交基 B，使 u 在基 B 中是对角的。

集 \(\mathrm{Sp}_{\mathrm{s}}(u)\) 是可数的，且当 E 无限时它不含 0（III, p. 90, 命题 5, \(b\) ）。对每个元素 \(\lambda \in \mathrm{Sp}_{\mathrm{s}}(u)\) ，用 \(\mathrm{N}_{\lambda}\) 表示 \(u - \lambda 1_{\mathrm{E}}\) 的幂零空间。它是有限维的（同处），且由于 \(u\) 正规，它与 \(u\) 的相对于 \(\lambda\) 的特征空间重合（EVT, V, p. 43, 命题 8 的推论）。诸空间 \(\mathrm{N}_{\lambda}\) 两两正交（I, p. 132, 第 5 目）。

设 F 是 E 的这样的子空间：它是诸空间 \(N_{\lambda}\) （对 \(\lambda \in \mathrm{Sp}_{\mathrm{s}}(u) - \{0\}\) ）的希尔伯特和。它是可数型的空间，在 u 下稳定，因为每个子空间 \(N_{\lambda}\) 在 u 下稳定。由于 \(N_{\lambda}\) 也是 \(u^{*}\) 的相对于 \(\overline{\lambda}\) 的特征空间（EVT, V, 同处），自同态 u 通过过渡到子空间诱导一个自同态 \(\widetilde{u}\) （\(F^{\circ}\) 的）。自同态 \(\widetilde{u}\) 是紧的（III, p. 5, 命题 3）且正规的（I, p. 135, 引理 4）。按构造，\(\widetilde{u}\) 的敏感谱包含于 \(\{0\}\) （事实上，\(\widetilde{u}\) 的每个特征向量也会是 u 的特征向量，从而若对应的特征值非零，它就属于某个空间 \(N_{\lambda}\) ）。于是 \(\widetilde{u}\) 的谱半径为零，从而 \(\widetilde{u}=0\) （因为 \(\widetilde{u}\) 正规；I, p. 108, 推论 1）。因此我们有 \(F^{\circ}\subset\operatorname{Ker}(u)\) 。

对每个 \(\lambda \in \mathrm{Sp}_{\mathrm{s}}(u)\) ，设 \(B_{\lambda}\) 是 \(N_{\lambda}\) 的一个标准正交基，设 \((e_j)_{j \in J}\) 是由诸 \(B_{\lambda}\) 的并组成的族；它是 F 的一个标准正交基。设 B 是 \((e_j)_{j \in J}\) 与 \(F^{\circ}\) 的一个标准正交基之并。这是 E 的一个标准正交基，且 \(u\) 在基 B 中是对角的。

推论 1. --- 设 u 是 E 的紧埃尔米特自同态。存在 E 的一个标准正交基 B，使 u 在基 B 中是对角的且其特征值为实数。

若 K = C，这立即由定理得出。假设 K = R。空间 E 是 \(\mathrm{E}_{(\mathbf{C})}\) 上的一个 R-结构（参见 A, II, p. 119）。自同态 \(u_{(\mathbf{C})}\) （\(\mathrm{E}_{(\mathbf{C})}\) 的）是紧的（III, p. 2, 注记 4）且埃尔米特的。设 \(B = (e_j)_{j \in J}\) 是 \(\mathrm{E}_{(\mathbf{C})}\) 的一个标准正交基，使得 \(u_{(\mathbf{C})} \in \mathcal{D}_{\mathrm{B}}(\mathrm{E}_{(\mathbf{C})})\) ，设 \(\lambda\) 是 \(u_{(\mathbf{C})}\) 的特征值族（定理 1）。我们有 \(\lambda \in R^{J}\) （I, p. 106, 命题 4）；由于线性映射 \(u_{(\mathbf{C})}\) 是 R-有理的，\(u_{(\mathbf{C})}\) 的相对于 \(\lambda_j\) 的特征空间对每个 \(j \in J\) 都是 R-有理的（参见 A, V, p. 60, 命题 6）。因此存在它的一个属于 E 的基，且更有理由存在 E 中的一个标准正交基。这些基之并是 E 的一个标准正交基 \(B_R\) ，使得 \(u \in \mathcal{D}_{\mathrm{B}_R}(\mathrm{E})\) 。

推论 2. --- 设 F 是希尔伯特空间，u 是从 E 到 F 的紧连续线性映射。存在 u 的始空间的一个标准正交基 \((e_{i})_{i\in\mathrm{I}}\) （\(\operatorname{Ker}(u)^{\circ}\) 的）、F 的一个标准正交族 \((f_{i})_{i\in\mathrm{I}}\) 以及一个族 \((\alpha_{i})_{i\in\mathrm{I}}\in(\mathbf{R}_{+}^{*})^{\mathrm{I}}\) ，使得 \(u(e_{i})=\alpha_{i}f_{i}\) （对所有 \(i\in I\) ）。

设 \(v = u^{*} \circ u\) 。这是 E 的一个紧的、正的（从而埃尔米特的）自同态。由推论 1，存在 E 的一个标准正交基 \((e_{j})_{j \in \mathrm{J}}\) ，使得 \(v\) 在此基中是对角的。它的特征值族 \((\lambda_{j})_{j \in \mathrm{J}}\) 包含于 \(\mathbf{R}_{+}^{\mathrm{J}}\) 。置 \(\alpha_{j} = \sqrt{\lambda_{j}}\) （对所有 \(j \in \mathrm{J}\) ）。设 I 是由 \(j \in \mathrm{J}\) （使 \(\alpha_{j} \neq 0\) ）组成的集。由于 \(v\) 紧，这是一个可数集。族 \((e_{i})_{i \in \mathrm{I}}\) 是 \(v\) 的始空间（即始空间 \(\operatorname{Ker}(u)^{\circ}\) ，\(u\) 的）的一个标准正交基（EVT, V, p. 43, 命题 8）。置 \(f_{i} = \frac{1}{\alpha_{i}} u(e_{i})\) （对 \(i\in \mathrm{I}\) ）。对 I 中任何 \(i\) 与 \(j\) ，我们有

\[
\langle f _ {i} | f _ {j} \rangle = \frac {1}{\alpha_ {i} \alpha_ {j}} \langle u (e _ {i}) | u (e _ {j}) \rangle = \frac {1}{\alpha_ {i} \alpha_ {j}} \langle v (e _ {i}) | e _ {j} \rangle = \frac {\lambda_ {i}}{\alpha_ {i} \alpha_ {j}} \langle e _ {i} | e _ {j} \rangle ,
\]

由此得出族 \((f_{i})_{i\in\mathrm{I}}\) 在 F 中是标准正交的。由于 \(u(e_{i})=\alpha_{i}f_{i}\) （对所有 \(i\in I\) ），推论得证。

定义 2. --- 沿用推论的记号，族 \((\alpha_{i})_{i\in I}\) 是 \(u\) 的相对于始空间的标准正交基 \((e_{i})_{i\in I}\) （该始空间为 \(u\) 的始空间）的奇异值族。

注记. --- 1) 该推论推广了 EVT, V, p. 54 的定理 2，后者对应于希尔伯特–施密特映射。

\begin{enumerate}
\def\labelenumi{\arabic{enumi})}
\setcounter{enumi}{1}
\tightlist
\item
  沿用推论的记号，我们有公式
\end{enumerate}

\[
u (x) = \sum_ {i \in I} \alpha_ {i} \langle e _ {i} | x \rangle f _ {i}\tag{1}
\]

对所有 \(x \in E\) 成立。

\begin{enumerate}
\def\labelenumi{\arabic{enumi})}
\setcounter{enumi}{2}
\tightlist
\item
  设 \(u\) 是 E 的紧正自同态。设 \(B = (f_i)_{i \in I}\) 是 E 的一个标准正交基，使得 \(u\) 在基 B 中是对角的（定理 1），设 \((\lambda_i)_{i \in I}\) 是 \(u\) 在此基中的特征值族。设 J 是由 \(i \in I\) （使 \(\lambda_i > 0\) ）组成的集；族 \((e_i)_{i \in J}\) 是空间 \(\text{Ker}(u)^{\circ}\) 的一个标准正交基。对每个 \(i \in J\) ，置 \(e_i = f_i\) 与 \(\alpha_i = \lambda_i\) 。我们得到 \(u(e_i) = \alpha_if_i\) ：族 \((\alpha_i)_{i \in J}\) 是 \(u\) 的相对于基 \((e_i)_{i \in J}\) 的奇异值族。
\end{enumerate}

